% !PRINTED 215
\weberpage{215}
%<<CONT
and the two fundamental forms become:
\weq{7}{\begin{aligned}
\Phi_1 &= x_1^4 + 2\sqrt{-3}\,x_1^2 x_2^2 + x_2^4 \\
\Phi_2 &= x_1^4 - 2\sqrt{-3}\,x_1^2 x_2^2 + x_2^4.
\end{aligned}}

Between the three fundamental forms $f, \Phi_1, \Phi_2$ a relation can be derived by eliminating $x_1, x_2$. From (7) one finds
\[
2(x_1^4+x_2^4) = \Phi_1 + \Phi_2
\]
\[
4\sqrt{-3}\,x_1^2 x_2^2 = \Phi_1 - \Phi_2,
\]
and from (1)
\[
f^2 = x_1^2 x_2^2 (x_1^4+x_2^4)^2 - 4x_1^6 x_2^6,
\]
from which one easily computes
\weq{8}{12\sqrt{-3}\,f^2 = \Phi_1^3 - \Phi_2^3.}

The substitutions $\Theta, \psi, \chi$, when represented with determinant
1, take the form:
\weq{9}{\begin{pmatrix} i, & 0 \\ 0, & -i \end{pmatrix},\quad
\begin{pmatrix} 0, & i \\ i, & 0 \end{pmatrix},\quad
\begin{pmatrix} \dfrac{1-i}{2}, & \dfrac{1-i}{2} \\[6pt] -\dfrac{1+i}{2}, & \dfrac{1+i}{2} \end{pmatrix},}
from which it follows, by the conditions established in \S.\,50, that
the corresponding group of orthogonal ternary substitutions
is real.

If one applies the substitutions (9) with the determinant 1
to the fundamental form $f(x_1, x_2)$, a very simple
\gesperrt{calculation shows that the fundamental form $f$ is an absolute invariant}
of the binary group $G$.

The \gesperrt{Hessian} determinant of $f$ also has this property:
\[
\begin{aligned}
H &= -\frac{1}{25}\left[f''(x_1,x_1)\,f''(x_2,x_2) - f''(x_1,x_2)^2\right] \\
&= (x_1^4+x_2^4)^2 + 12\,x_1^4 x_2^4,
\end{aligned}
\]
which is nothing other than the product of the two functions
$\Phi_1, \Phi_2$; while the functions $\Phi_1$ and $\Phi_2$ themselves, under the third
of the substitutions (9), acquire a third root of unity as a
factor.
